% Standalone problem document, generated from the adjacent README.md.
% Compile with XeLaTeX. Mathematical content is maintained in Markdown.
\documentclass[11pt,a4paper]{article}
\usepackage[fontsize=10.5pt]{scrextend}
\usepackage[margin=22mm,headheight=15pt,headsep=9mm,footskip=11mm]{geometry}
\usepackage{fontspec}
\setmainfont{texgyrepagella}[Extension=.otf,UprightFont=*-regular,BoldFont=*-bold,ItalicFont=*-italic,BoldItalicFont=*-bolditalic]
\setsansfont{texgyreheros}[Extension=.otf,UprightFont=*-regular,BoldFont=*-bold,ItalicFont=*-italic,BoldItalicFont=*-bolditalic]
\setmonofont{texgyrecursor}[Extension=.otf,UprightFont=*-regular,BoldFont=*-bold,ItalicFont=*-italic,BoldItalicFont=*-bolditalic,Scale=MatchLowercase]
\usepackage{amsmath,amssymb}
\usepackage{unicode-math}
\setmathfont{texgyrepagella-math.otf}
\usepackage{xcolor,microtype,enumitem,needspace,fancyhdr}
\usepackage{bookmark}
\definecolor{ink}{HTML}{17344B}
\definecolor{accent}{HTML}{197D80}
\definecolor{muted}{HTML}{596773}
\hypersetup{colorlinks=true,linkcolor=accent,urlcolor=accent,pdftitle={TR-20 - Rayleigh--Ritz
discriminant degrees for rank-one
matrices},pdfauthor={Open Problems in Numerical Linear Algebra}}
\urlstyle{same}
\setlength{\parindent}{0pt}
\setlength{\parskip}{5pt plus 1pt minus 1pt}
\linespread{1.04}
\setlength{\emergencystretch}{3em}
\widowpenalty=10000
\clubpenalty=10000
\displaywidowpenalty=10000
\allowdisplaybreaks[1]
\setlist{nosep,leftmargin=1.5em}
\providecommand{\tightlist}{\setlength{\itemsep}{0pt}\setlength{\parskip}{0pt}}
\providecommand{\pandocbounded}[1]{#1}
\pagestyle{fancy}
\fancyhf{}
\fancyhead[L]{\sffamily\footnotesize\color{muted}OPEN PROBLEMS IN NUMERICAL LINEAR ALGEBRA}
\fancyhead[R]{\sffamily\bfseries\footnotesize\color{accent}TR-20}
\fancyfoot[L]{\parbox[b]{0.9\textwidth}{\sffamily\fontsize{8}{10}\selectfont\color{muted}Editorial ratings; a bounded search cannot certify that a problem remains unsolved.\\Literature check: 2026-09-11}}
\fancyfoot[R]{\sffamily\footnotesize\color{muted}\thepage}
\renewcommand{\headrulewidth}{0.3pt}
\renewcommand{\headrule}{\hbox to\headwidth{\color{accent}\leaders\hrule height \headrulewidth\hfill}}
\makeatletter
\renewcommand{\section}{\Needspace{5\baselineskip}\@startsection{section}{1}{0pt}{12pt plus 2pt minus 2pt}{4pt}{\sffamily\bfseries\large\color{ink}}}
\renewcommand{\subsection}{\Needspace{4\baselineskip}\@startsection{subsection}{2}{0pt}{9pt plus 1pt minus 1pt}{3pt}{\sffamily\bfseries\normalsize\color{ink}}}
\makeatother
\setcounter{secnumdepth}{0}
\begin{document}
{\sffamily\small\color{muted}Tensor computations\par}
\vspace{3pt}
{\raggedright\hyphenpenalty=10000\exhyphenpenalty=10000\sffamily\bfseries\fontsize{21}{24}\selectfont\color{ink}Rayleigh--Ritz
discriminant degrees for rank-one matrices\par}
\vspace{7pt}
{\sffamily\small\textbf{Difficulty:} challenging\quad\textbf{Importance:} interesting
to specialist\par}
{\sffamily\small\bfseries\color{accent}Status: Solved\par}
\vspace{4pt}
{\color{accent}\hrule height 0.8pt}
\vspace{6pt}
\textbf{Rating rationale:} Challenging because general
discriminant-degree formulas must account for degeneracies and the
excluded isotropic locus; specialist importance concerns algebraic
conditioning of rank-one Rayleigh--Ritz optimization.\\
\textbf{Source:} Borovik--Friedman--Hoşten--Pfeffer, Conjecture 3.18 in
v2.

\subsection{Resolution --- 2026-09-11}\label{resolution-2026-09-11}

\textbf{Affirmative resolution by Matthew J. Colbrook}, DAMTP,
University of Cambridge. Theorem 1 proves both displayed reduced-degree
formulas for every \(n\ge2\) in the exact complex bilinear Segre model,
including boundary and multiplicity analysis.

\href{https://github.com/ajt60gaibb/OpenProblemsInNLA/blob/main/references/colbrook-recovered-tensors-2026-09-11/manuscripts/TR-20.pdf}{Complete
proof, Theorem 1 and §§2--8} ·
\href{https://github.com/ajt60gaibb/OpenProblemsInNLA/blob/main/references/colbrook-recovered-tensors-2026-09-11/manuscripts/TR-20.tex}{TeX}
·
\href{https://github.com/ajt60gaibb/OpenProblemsInNLA/blob/main/references/colbrook-recovered-tensors-2026-09-11/verification/reviews/TR-20-review.md}{Independent
PASS review} ·
\href{https://github.com/ajt60gaibb/OpenProblemsInNLA/blob/main/references/colbrook-recovered-tensors-2026-09-11/README.md}{Authorship
and provenance}.

The AI-assisted proof passed independent agent review, not external
human peer review or formal certification. Ratings are historical; no
priority claim is made.

\subsection{Problem statement}\label{problem-statement}

For integers \(m,n\ge2\), put \(N=mn\) and let

\[
V_{m,n}=\{[\operatorname{vec}(ab^T)]:
0\ne a\in\mathbb C^m,\ 0\ne b\in\mathbb C^n\}
\subset\mathbb P^{N-1}.
\]

This is the Segre variety of nonzero rank-one \(m\times n\) matrices
modulo scaling, in the ordinary entry coordinates. For a complex
symmetric matrix \(H\in\mathbb C^{N\times N}\) define

\[
R_H([\psi])=\frac{\psi^TH\psi}{\psi^T\psi},
\qquad
[\psi]\in V_{m,n},\quad \psi^T\psi\ne0.
\]

Transpose here is bilinear transpose, not conjugate transpose. A
critical point is one at which the differential of \(R_H\) restricted to
the smooth variety \(V_{m,n}\) vanishes. It is degenerate if its Hessian
in local coordinates on \(V_{m,n}\) is singular; at a critical point
this condition is independent of the local coordinates.

Let \(\Delta^{\mathrm{ni}}_{m,n}\) be the Zariski closure in
\(\mathbb P(\operatorname{Sym}^2\mathbb C^N)\) of the set of nonzero
symmetric matrices \(H\), modulo scaling, for which such a degenerate
critical point exists with \(\psi^T\psi\ne0\). This is the source's
nonisotropic discriminant.

\textbf{Conjecture.} For every integer \(n\ge2\), these discriminants
are hypersurfaces and their degrees are\\
\[
\deg\Delta^{\mathrm{ni}}_{2,n}
=24\binom{n+1}{3},
\qquad
\deg\Delta^{\mathrm{ni}}_{3,n}
=24n^2\binom n2.
\]

Degree means projective degree of the reduced hypersurface, equivalently
its intersection number with a general projective line. Both formulas
are retained as one source conjecture.

\subsection{Why it matters in numerical linear
algebra}\label{why-it-matters-in-numerical-linear-algebra}

Minimizing \(R_H\) over real rank-one matrix states is a constrained
energy problem underlying tensor variational algorithms. The
discriminant records degeneracies of stationary states, where local
solution branches can meet and conditioning can deteriorate.

\newpage
\subsection{References}\label{references}

\begin{enumerate}
\def\labelenumi{\arabic{enumi}.}
\tightlist
\item
  Viktoriia Borovik, Hannah Friedman, Serkan Hoşten, and Max Pfeffer,
  \href{https://arxiv.org/html/2512.06939v2}{\emph{Numerical Algebraic
  Geometry for Energy Computations on Tensor Train Varieties}},
  arXiv:2512.06939v2 (2026), §3.2, Proposition 3.7, definition before
  Example 3.11, Conjecture 3.18; §7.2, Table 4.
\item
  Flavio Salizzoni, Luca Sodomaco, and Julian Weigert,
  \href{https://arxiv.org/html/2510.17760v1}{\emph{Nonlinear Rayleigh
  quotient optimization}}, 2025, §§2 and 5.
\end{enumerate}

\subsection{Status check --- 2026-09-10}\label{status-check-2026-09-10}

Rechecked \href{https://arxiv.org/html/2512.06939v2}{Borovik et al.~v2,
Conjecture 3.18 and Table 4}, and searched the title, authors,
discriminant and conjecture number. The June 2026 revision retains both
formulas as one conjecture. Small-dimensional numerical evidence and
generic critical-point counts are not proofs of the stated general
discriminant degrees. No later resolution was located.
\end{document}
